\clearpage\SourceLeaf{22}
\begin{center}
\bfseries ZUR THEORIE DER CUBISCHEN RESTE.
\end{center}



\begin{center}
[1.]
\end{center}

\[ p=mm+3nn \text{ eine Primzahl.} \]

\noindent\hspace*{\parindent}Man setze für irgend eine ganze durch $p$ nicht theilbare Zahl $x$
\[
\left[\dfrac{2mx}{p}\right]=\alpha,\qquad \left[\dfrac{(3n-m)x}{p}\right]=\beta,\qquad \left[\dfrac{(-3n-m)x}{p}\right]=\gamma.
\]
\noindent Hieraus wird, wie man leicht sieht,
\[
\dfrac{2mx}{p}-\alpha+\dfrac{(3n-m)x}{p}-\beta+\dfrac{(-3n-m)x}{p}-\gamma
\]
die Summe aus drei echten Brüchen, und muss daher exclusive zwischen den
Grenzen 0 und 3 liegen; da aber offenbar dieselbe Summe $=-\alpha-\beta-\gamma$, also
eine ganze Zahl ist, so kann sie keinen andern Werth als $+1$ und $+2$ haben.
Es ist also entweder
\[ \alpha+\beta+\gamma=-1\quad\text{oder}\quad \alpha+\beta+\gamma=-2. \]

\noindent Hieraus folgt ferner, dass die Grössen $\alpha, \beta, \gamma$ nicht alle drei unter einander
nach dem Modulus 3 congruent sein können; unter den Differenzen $\alpha-\beta$, $\beta-\gamma$,
$\gamma-\alpha$ wird also wenigstens Eine (und eben deswegen wenigstens noch eine,
